% % 春休みのひととき　\\
\subsection{春假的片刻时光}

% 休みの日に千歳が周の家に遊びに来る事は、珍しくない。そこに真昼が居るからであって周がおまけなのは目に見えているが、特に用事もないので場所を提供するくらいはいいだろう、という事で千歳を受け入れる事が多かった。\\
休息日千岁来周家玩，这并不罕见。明眼人都看得出来，这是因为真昼在这里，周不过是个附带的。但反正也没什么特别的事，提供个场地也无妨，所以周大多时候都会接纳千岁。\\

% 「うーん春休みという学生幸せ期間の一つなのにあんまり休みを感じない」
「唔～明明春假是学生幸福时期之一，却没什么放假的感觉」

% 「そりゃ今課題真っ只中だからな」\\
「那是因为现在正被作业追着跑嘛」\\

% 　今日は意気揚々とやる気なくという矛盾したテンションで千歳だけが来訪している。
今天只有千岁一个人来访，情绪处于一种既意气风发又毫无干劲的矛盾状态中。

% 　リビングで課題に向き合っている千歳は、やはり中々問題を解くスピードも上がらない。長期休暇という学生にとっては待ち侘びた時間だろうが、悲しいかなこの春から受験生、遊び一辺倒にするのは不可能である。\\
千岁在客厅里对着作业奋斗着，解题速度依然提不上来。长假对学生来说本该是翘首以盼的时光，但可悲的是，从这个春天开始就要成为备考生了，想把时间全用来玩是不可能的。\\

% 　ローテーブルに突っ伏しはしないものの肘をついて気を進まなそうにシャーペンを回している千歳だが、これでも彼女から言い出した事なのでやる気はある方だ。\\
千岁虽然没有趴在矮桌上，但也撑着手肘，一副提不起劲的样子转着自动铅笔。不过这也是她自己提出来的，所以还算是比较有干劲。\\

% 「うえーん何で学年変わるこの休みまで課題あるのぉ」
「呜，为什么到了换学年的这个假期还有作业啊」

% 「受験生になるんだからそりゃ多くの時間を勉強に割くためだろ」
「既然要成为备考生了，那当然是为了把更多的时间分配给学习」

% 「そういう事を言ってるんじゃないの。ナンセンスだねえ」
「我说的不是这个意思。真没品位呢」

% 「俺の事を何と言おうと構わないが課題はその分解くのが遅くなるぞ」
「随你怎么说我都没关系，但那样的话你做题的速度会更慢哦」

% 「周が意地悪言う。しかも余裕こいてるし」
「周真坏。而且还一副游刃有余的样子」

% 「先に終わらせて毎日予習復習してますので」
「因为我已经提前做完，现在每天都在预习和复习了」

% 「ぐあー！　優等生が！」\\
「呜哇——！你这个优等生！」\\

% 　千歳とローテーブルを挟んで周もノートを開いているが、これは参考書を解いているだけで課題は既に片付いて通学鞄の中。\\
隔着矮桌，周也摊开了笔记本。不过这只是在教辅书上的题而已，作业早就搞定塞进书包里了。\\

% 　バイトで忙しくなるのは見えていたので、初日にさっさと終わらせて後は隙間時間に予習復習と詰め込んでいるのだが、これだけで優等生になれるなら受験生の多くは優等生である。\\
因为预见到之后打工会很忙，所以在放假第一天就赶紧把作业做完，剩下的空闲时间都塞满了预习和复习。不过如果仅凭这样就能成为优等生的话，那大部分备考生都是优等生了。\\

% 　余裕持ったスケジュール管理を、そして反復作業を、と両親どころか真昼にも言われているので、流石の周もオンオフはきっちりして休みと勉強、バイトを併行させている。\\
连父母甚至真昼都叮嘱过要进行有余量的日程管理以及反复练习，所以就连周也是把状态切换得很清楚，将休息、学习和打工并行安排。\\

% 「文句言いつつ殆ど終わってますよね千歳さん」\\
「千岁虽然嘴上抱怨，但也差不多快做完了吧」\\

% 　うだうだしている千歳ではあるが、真昼の言う通り、課題の進み具合を見るに家できっちりやっているようだ。\\
虽然千岁看起来慢吞吞的，但正如真昼所说，从作业的进度来看，她在家里也有好好在做。\\

% 「……そりゃあ？　ちゃんと頑張るって決めてるし、今までと同じじゃ駄目じゃん」
「……那是当然啦？我都决定要好好努力了，总不能和以前一样吧」

% 「おお、千歳が成長した……」
「哇哦，千岁成长了……」

% 「馬鹿にしてない？」
「你是不是在小看我？」

% 「今までのお前の言動と実行したものを遡ってみろ」
「你回顾一下你以前的言行和实际做到的事情」

% 「……今頑張ってるから！」
「……我现在可是很努力的！」

% 「まあそれはよい事で。……樹はよかったのか？」
「这倒是好事。……树没关系吗？」

% 「いっくん今日は塾だってさ。お母さんに頭を下げて入れてもらったって」
「阿树说今天要去补习班。说是拜托妈妈才让他进的」

% 「大輝さんに頭下げない辺りに微かな抵抗を感じる」\\
「尽量不向大辉叔叔低头这一点，能感觉到他那微弱的抗拒呢」\\

% 　受験生という事もあって三年生から塾や予備校に通い出す生徒も珍しくはない。樹もその例に漏れなかった。
因为是备考生，所以从高三开始去补习班或预备校的学生也不少见。树也不例外。

% 　父親になるべく言いたくないのは彼らしいが、もう少し歩み寄っておいても損はないんだがなあ、とも思ってしまうあたり、やはり当事者と見えているものが違うのだろう。\\
不想向父亲低头这确实很符合他的性格，但周还是觉得，要是两人能走近一些对彼此都没坏处。果然当事人和旁观者看到的东西是不一样的。\\

% 「まあ、大輝さんが知らない間にとんでもなく賢くなって見返してやるって言ってたからね。私も行かなきゃなんだけどねー、お兄ちゃんの大学費用と私の大学費用で結構カツカツになってるらしいし、塾やとか予備校一年丸々は厳しいっぽいからもうちょっと後からだなー。周達は？」
「他说，要在大辉叔叔不知道的时候变得超级聪明，然后让他刮目相看呢。其实我也必须得去啦～不过据说哥哥和我的大学费用加起来家里手头已经挺紧的了，整整一年都去补习班或预备校好像挺困难的，所以打算稍微晚一点再去～。周你们呢？」

% 「俺はバイト終わるまでは時間の都合つけにくいしオンラインと映像授業の塾にして夏からは予備校。五月とか七月の模試の結果で両親と相談してその辺りはもうちょい細かく詰めるつもり」\\
「我打工结束前时间不好安排，所以先选了线上和视频授课的补习班，夏天开始再去预备校。打算根据五月或者七月的模拟成绩和父母商量，到时候再更详细地规划一下」\\

% 　流石に周の自認している学力では学校の勉強だけで確実に志望校に受かるとは言い切れないので、高校三年生丸々受験に向けて取り組む人達程ではないにせよ受験生として堂々と名乗れる程度には勉強するつもりだ。
以周自己认知的学习能力，确实不敢断言光靠学校的学习就能稳上志愿学校，所以虽然不像那些高三整整一年都在为考试拼命的人那样，但他也打算学到能堂堂正正自称为备考生的程度。

% 　両親も周が勉学に励む事自体は賛成らしく、応援は惜しまない上金銭も心配しなくていいとの仰せで有り難い限りである。\\
父母对周努力学习这件事本身似乎也很赞成，不仅不吝惜支持，还说不用担心钱的问题，这真是再感激不过了。\\

% 「どっちも自分で自分を律するから出来るやつじゃん」
「你们两个都是那种能自律所以才能做到的家伙嘛」

% 「受験生が受験シーズンに自分を律しないでいつ律するんだ」
「备考生在备考季不自律，要等到什么时候才自律啊」

% 「受験本番？」
「正式考试的时候？」

% 「それは諦めか投げやりって言うんだぞ……」
「那叫放弃或者破罐子破摔哦……」

% 　受験本番だけ真面目にやるのは最早受験を投げていると同義だろう。
只有在正式考试时才认真，这跟放弃考试是一个意思了吧。

% 「まひるんは？」
「昼儿呢？」

% 「私も周くんと似たようなものですね。幸いというのか、受験に必要な授業は学び終えていますし、後は受ける大学に合わせて勉強内容を変えて対策していこうかなと。周くんと学部は違えど大学は一緒ですからその辺り周くんとも連携取れたらいいですよね」
「我和周君差不多吧。不知道算不算幸运，考试需要的课程我已经学完了，之后打算根据要报考的大学，针对性地改变学习内容。虽然我和周君的学院不同，但大学是一样的，如果能在这些方面和周君互相协助就好了」

% 「……周りが真面目だと私のいい加減さが浮き彫りになって辛い」
「……周围的人都这么认真，更凸显了我的随便，好难受」

% 「千歳さんも頑張っていると思いますよ。少なくとも前からしたら見違える程進歩しています」
「我觉得千岁也很努力了哦。至少和以前相比，进步可以说是让人刮目相看」

% 「だめまひるん、私を甘やかさないで……！」
「不行昼儿，不要惯着我……！」

% 「じゃあおやつのケーキは要らないと」
「意思是说不需要蛋糕来当点心了」

% 「それは話が別！」\\
「那是两码事！」\\

% 　やはり課題を暫くこなして糖分足りなくなっていたな、という周の予想通りケーキに食い付くので、周も真昼も笑ってキッチンに向かう。\\
果然是做了一会儿作业糖分不足了吧，就像周预料的那样，一提到蛋糕她就上钩了。周和真昼都笑了笑，走向厨房。\\

% 　周がケーキを皿に移している間に真昼が紅茶の用意をするという別に何もおかしくない光景なのに「ひゅーひゅー」と囃し立てるような飛んでくる。何やってんだあいつ、とは思うが反応したらそれこそ千歳の思うがままなので無視しておいた。\\
周把蛋糕盛到盘子里的这段时间，真昼准备红茶，这明明是再正常不过的光景，客厅里却飞来了咻咻的起哄声。这是在干什么啊，周心想。但要是回应了那正好如了千岁的意，所以周选择无视。\\

% 　真昼ももう慣れているのか千歳の揶揄には一々反応しておらず、リビングで騒ぐ人間が一人という状況になっている。\\
真昼似乎也习惯了，对千岁的调侃并没有一一回应，客厅里呈现出只有一个人在喧闹的状况。\\

% 「真昼、俺今回ミルクありがいい」
「真昼，我这次想加点牛奶」

% 「じゃあミルクピッチャーに多めに入れないとですね。あ、周くんそっちの皿よりこっちの方がケーキのサイズ的に丁度いいですしカップともお揃いですよ」
「那得在奶盅里多倒一点呢。啊，周君，用我这边的盘子，比你那的更适合蛋糕的尺寸，而且和杯子也是一套的」

% 「こっちね。揃えてよかったな」
「这个对吧。还好买了成套的」

% 「お二人さん私の存在目に入ってる？」
「你们两位，眼里有我这个人的存在吗？」

% 「視界には入ってないな」
「没在视野里呢」

% 「声はお元気そうだなと」
「听声音倒是觉得挺有精神的」

% 「……やっぱりいちゃついてるじゃーん」
「……果然还是在调情嘛～」

% 「はいはい」\\
「是是是」\\

% 　結局話はそこに帰結するらしい千歳に苦笑いしながら、今日千歳のために用意してあったケーキを皿に載せて真昼の用意した紅茶と共にトレイでリビングに持っていく。\\
最终话题似乎总会绕到这上面，周对千岁回以苦笑，将今天为千岁准备的蛋糕盛在盘子里，和真昼准备的红茶一起放在托盘上端到了客厅。\\

% 　課題に飽きたらしい千歳は文字通り諸手を挙げて休憩に大賛成の姿勢を見せており、千歳お気に入りのパティスリーの季節限定ケーキを見て目を輝かせていた。\\
千岁似乎做作业做烦了，高举双手赞成休息，看到她最喜欢的甜品店的季节限定蛋糕，眼睛都亮了起来。\\

% 「やっぱやだなーって思うのが三年生からみんなピリピリしだす事なんだよね」\\
「果然还是觉得很讨厌啊——从高三开始大家就都变得神经紧绷起来了」\\

% 　切り替えが早い千歳はもぐもぐと実に美味しそうにケーキを切り分けて口にしているが、話しているのは甘くなさそうな現実の話だ。\\
千岁转换心情很快，她嘴里嚼着美味的蛋糕，说出的却是毫不甜蜜的现实话题。\\

% 　もう三年生に進級する事になるが、やはり二年生の時とは雰囲気が異なってくるだろう。何せこれから一年は重要なものになる。今までのような気楽な雰囲気とは違ってくるだろう。\\
马上就要升入高三了，气氛肯定会和高二时有所差异。毕竟接下来的一年将是非常重要的一年，和迄今为止那种轻松的氛围截然不同。\\

% 「程よい緊張感ならいいんだけどな。おいそれと他人の成績には触れない方が無難な空気になるだろうな。元々触れるつもりもないけど」
「如果有适度的紧张感倒还好，大概就是会弄得不能轻易去提及别人的成绩吧。我倒是本来也不会这么干」

% 「比べちゃうもんねえみんな」
「大家都会忍不住去比较的嘛」

% 「相対的評価で自分の位置を確認しないと不安なんだろうなあと」
「如果不通过相对评价来确认自己的位置，就会感到不安吧」

% 「うーん憂鬱」
「呜～好郁闷」

% 「こればっかりは受験生ならよくある事なのではないかと。比較の負のループに陥らないように気を付けていきたいですね」\\
「这种事对备考生来说是很常见的吧。希望我们能注意不要陷入互相对比的恶性循环里去」\\

% 　穏やかな笑顔で品よく紅茶を飲んでいる真昼は、全体的に余裕があるからか進級にも臆した様子はなさそうだ。どんと構える姿勢は頼もしいの一言に尽きる。\\
真昼带着温和的笑容优雅地喝着红茶。或许是因为整体上游刃有余，看起来对升学并没有感到畏怯。那份从容不迫的姿态，只能用可靠来形容。\\

% 「うーんやっぱ暗い話題やめとこ」
「唔～还是不要聊这种沉重的话题了吧」

% 「千歳が言い出したんだけどな……」
「明明是你先挑起的话题……」

% 「まあまあ。明るい話題明るい話題！　もうすぐ春休み終わっちゃうけど、二人は何かしたい事とかないの？」
「好了好了。聊点开心的话题、开心的话题！春假马上就要结束了，你们俩有没有什么想做的事？」

% 「春休みが終わる事を明るい話題……？」
「春假要结束了算开心的话题……？」

% 「おだまり。ほら周は何かないの？」
「闭嘴。周有没有什么想做的？」

% 「したい事、なあ」\\
「想做的事啊……」\\

% 　したい事、と急に言われても中々に思いつかない。するべき事は幾らでもあるが、したい事、となると途端に候補すら思い浮かばない。\\
突然被问到想做的事，还真是一时想不起来。应该做的事倒是有很多，但要说想做的事，就连个候补选项都想不出来。\\

% 「課題はした、新しい参考書も買った。春の大掃除はした、布団をクリーニングに出した、制服もクリーニングに出した、靴も洗濯と手入れはしたし定期便も確認したし……」
「作业做完了，新的教辅书也买了。春季大扫除做完了，被子送去干洗了，制服也送去干洗了，鞋子洗了也保养了，定期配送的东西也确认过了……」

% 「やりたい事じゃなくてやらなきゃいけない事だし途中から所帯じみてるんですよお兄さん」
「这不是想做的事，而是必须做的事吧，而且听起来越来越像个家庭主夫了啊小哥」

% 「こちとら一人暮らしだぞ」\\
「我可是自己一个人住的」\\

% 　これはどうしようもない事ではあるが、一人暮らしは家事の負担が全部自分にかかってくる。周も自分の負担を真昼に押し付ける訳にはいかないので、当然ではあるが自分の事は自分できっちりこなすようにしている。
这虽然也是无可奈何的事，但一个人住所有的家务负担都会落在自己身上。周也不可能把自己的负担推给真昼，所以理所当然地，他会尽量把自己的事情都妥善处理好。

% 　こういう時に実家だと大体してもらえる事が多いので、一人暮らしが始まって両親の偉大さに感謝した事は今でも覚えている。\\
这种时候如果是在老家，大部分事情父母都会帮忙做好。他至今依然记得开始一个人住的时候，有多么感谢父母的伟大。\\

% 「一人、ふーん」
「一个人，哼～」

% 「何だよその顔」
「你那是什么表情」

% 「何でもないですー」\\
「没什么啦～」\\

% 　何か言いたそうだが何も言わないつもりの千歳にやや苛立ちながらもどうせ意味のない事かからかいのどちらかだろうと納得して、周はそれ以上追及せずに自分の分のプリンを口に放り込む。\\
周对这副似乎想说什么但又不打算说的千岁感到有些烦躁，但他也认定这不是无意义的话就是调侃，于是没有继续追问，而是把自己的那份布丁放进了嘴里。\\

% 「シケてる周は置いておくとして、まひるんは？」
「扫兴的周先不管了，昼儿呢？」

% 「……そうですねえ、春休みにしたい事、ですか」\\
「……这个嘛，春假想做的事吗」\\

% 　真昼もこの春休みはあまり外に出かけておらず、どちらかといえば春休み始めにあった出来事から心を落ちつけるために大人しくしていた印象がある。
真昼这个春假也没怎么出门，给人的印象是为了平复春假刚开始时发生的那件事带来的情绪，一直安分地待在家里。

% 　もし真昼がなにかしたいという事であれば周も協力するのだが、真昼の反応的には周と同じように今特段したい事はなさそうな気配が漂ってきていた。\\
如果真昼有什么想做的事，周也会配合她。但从真昼的反应来看，似乎和周一样，现在并没有什么特别想做的事。\\

% 「周といちゃいちゃするとかでもいいんだよ？」
「比如和周亲亲热热之类的也可以哦？」

% 「それは特別にしたい事ではないですし」
「那并不是什么特别的事」

% 「えっ」
「哎」

% 「違います！　春休みに、特別に、したいとかではなくて、その、出来ればそういうのは普段からという意味で……！　勿論春休みを利用した過ごし方もしたいですけど！」
「不是！我的意思是，并不是只有在春假才特别想做，那个，如果可以的话，我希望平时就能这样……！当然我也想利用春假好好度过啦！」

% 「分かってる分かってる、真昼にとって俺と過ごすのはいつもになったんだなって」
「我懂我懂，对真昼来说，和我一起度过的时光已经成为日常了对吧」

% 「ううう」
「呜呜呜」

% 「うーんいちゃついてますねえ。ノルマ達成？」
「唔～在调情呢。达成指标了？」

% 「ノルマなんてもの設けなくても千歳が居ない所で勝手にいちゃつきます」
「哪用得着什么指标，在千岁不在的地方我们也会自己调情的」

% 「のけものだー」
「我被排挤了～」

% 「逆に人前でする度胸あんのか……いやあったわこいつら」\\
「反过来说，你们有在别人面前调情的胆量吗……不对，有呢，这家伙们」\\

% 　千歳はよくこちらの事を言ってくるが、千歳達も大概である。寧ろ人目がある所でいちゃつくのは圧倒的に彼女達であり、周と真昼は親しい人間しか居ない場所でしかそういったやり取りはしない。していないと、思う。\\
千岁虽然经常说周的事，但千岁他们自己也差不多。倒不如说，更经常在有人的地方调情的是他们俩。除非周围全是亲密的人，否则周和真昼是不会有那样的互动的。应该，没有。\\

% 「私達だって周が居なくてもいちゃつくし周が居てもいちゃつくもーん」
「我们也是就算周不在也会调情，周在也会调情嘛～」

% 「元祖バカップルの名は伊達じゃないな」
「元祖笨蛋情侣的名号果然名不虚传」

% 「それ勝手に周が言ってるだけなんだけど。まあ否定はしないけど」
「那只是周自己随便说说的啦。虽然我也不否认就是了」

% 「しないのかよ」
「不否认啊」

% 「それたけラブラブという事の証ですからー」\\
「那是证明我们非常恩爱的证据嘛～」\\

% 　ふふん、と胸を貼る千歳であるが、一時期曇りに曇っていたのでこうして元気に樹と仲睦まじい様子を見せてくれるのは安堵の材料になるし、二人の精神衛生を測る上で重要だったりする。
哼哼，千岁挺起了胸膛。有一段时间她情绪非常低落，所以像这样充满活力地展现出和树和睦相处的样子，让人感到安心，这也是衡量他们俩心理健康的重要指标。

% 　年末にかけて大喧嘩した後から今まで、取り敢えず揉め事もなく日々を過ごしているので、心配ないだろう。\\
从年末那次大吵一架之后到现在，他们没什么纠纷地过着日子，应该不用担心了吧。\\

% 　樹も塾やバイトがあるそうなのですれ違いが起きないといいなとは思うのだが、自分にも降り掛かってきそうな言葉なので口を噤んでおく。\\
树好像也有补习班和打工，希望他们不要因此产生错位和隔阂就好了，不过这话感觉也会落到自己头上，所以周还是闭上了嘴。\\

% 「今日は樹が居なくて残念だったな」
「今天树不在真遗憾啊」

% 「ほんとそれ。こうなったら私もまひるんといちゃつきます」
「真的。既然这样，那我也和昼儿亲亲热热」

% 「お好きにどうぞ」
「请自便」

% 「まひるん売られちゃったよ？」
「昼儿被卖掉了哦？」

% 「人聞きの悪い事言うな。二人の友情に水を差す程野暮じゃないし、最終的には絶対俺の所に帰ってくるんだから」\\
「别说得那么难听。我还没不识趣到去给你们俩的友情泼冷水，而且最终她绝对会回到我这里的」\\

% 　千歳がおふざけでやっているのは百も承知であるし、そもそも千歳と真昼のソレは友情の範囲から逸脱するものではない。\\
周很清楚千岁是在开玩笑，而且千岁和真昼的那种互动并没有超出友情的范畴。\\

% 　真昼も楽しそうにしているのだから周が止める事はまずないし、そもそも千歳といちゃついた所で、真昼は周の所に戻ってくるのは分かっている。どれだけ千歳といちゃつこうが、真昼は自分の事を選ぶ自信があった。そもそも千歳は千歳で樹の所に帰るだろうに。\\
既然真昼看起来也很开心，周当然不会去阻止。说到底，就算她和千岁调情，真昼也会回到周身边，这一点他很清楚。无论怎么和千岁调情，周都有自信真昼会选择自己。再说，千岁也要回的树的身边嘛。\\

% 　だから存分にいちゃついてもらって、とミルクたっぷりのミルクティーを飲んで精神的には少し離れた位置から真昼と千歳を眺める周に、二人は顔を見合わせた。相変わらず仲良しそうだった。\\
所以就让她们尽情调情吧。周喝着加浓了牛奶的奶茶，在精神上稍微拉开了一点距离看着真昼和千岁。两人相视一笑，看起来依然很要好。\\

% 「うわー、うわー」
「哇～哇～」

% 「なんだよその顔」
「你那是什么表情」

% 「これが自信を得た周なんだなあと」
「这就是有了自信的周吗」

% 「俺の事馬鹿にしてます？」
「你是不是在小瞧我？」

% 「してませーん」\\
「没有哦～」\\

% 　していないと言う割に顔がどう考えてもしているのだが、真昼まで何故か納得したような顔だった。\\
虽然嘴上说着没有，但从表情怎么看都并非如此。不知为何连真昼也露出了一副心领神会的表情。\\

% 「よかったねえまひるん、昔のへたれな周じゃなくなって」
「太好了呢昼儿，以前那个懦弱的周已经不在了」

% 「……それはそれで困る時があるのですけど」
「……有时那也挺让人头疼的」

% 「えっ何々ききたーい」
「哎？什么什么，我想听～」

% 「おいやめろ。真昼もだめ」
「喂别闹了。真昼也不准说」

% 「えー」
「唉——」

% 「けちー」
「小气～」

% 「仲いいなほんと……」
「关系真好啊……」

% 「そりゃあ仲良しですので」
「那是自然啦」

% 「あっそ」
「啊是吗」

% 「そこでヤキモチやいてるくらい言ったらいいのにー」
「这种时候你吃点醋也没事的哦～」

% 「何で妬くんだよ……お前は俺が樹と仲良くしてたら妬くのか？」
「吃什么醋啊……要是我和树关系很好，你会吃醋吗？」

% 「えっいっくんとそういう関係が」
「哎？你和阿树是那种关系？」

% 「何でそうなるんだよ！」
「为什么会变成那样啊！」

% 「ムキになっちゃってー」
「还较真了呢——」

% 「お前は軽率に自分の彼氏にあらぬ疑惑を吹っ掛けるのはやめろ。草葉の陰で泣いてるぞ」
「你别轻易对自己的男朋友抛出莫名其妙的嫌疑啊。他会死不瞑目的哦」

% 「勝手に殺さないで下さーい。忙殺されてるかもしれないけどー」\\
「请不要随便杀人～虽然他可能忙的要死了～」\\

% 　ここには居ない樹は、今きっと同じ志を持った人間と机を並べてペンを握り締めている事だろう。何だかんだ気質的にはかなり真面目な樹らしくもあり、大輝と向き合う事を決めたからこその姿勢だ。\\
不在这里的树，现在肯定正和有着同样志向的人并排坐在桌前，紧握着笔吧。不管怎么说，从气质上来说这也很符合相当认真的树，这也是他决定面对大辉后所展现出的姿态。\\

% 　本人にはあまり言えないが、立派だと思っているし、尊敬する所でもある。\\
虽然不太好对本人说，但周觉得很了不起，也是他尊敬的地方。\\

% 「あいつも忙しそうだなあ」
「那家伙看起来也很忙啊」

% 「だね」\\
「是啊」\\

% 　少ししんみりしてしまったが、樹もずっと勉強しっぱなしで会えないという訳ではない。学校で会うだろうし、用事を作ればいつだって顔くらい見られる。\\
虽然气氛稍微变得有些伤感，但树也不是为了学习忙到没法见面。在学校里也会见面，找点事做的话，随时见个面总是没问题的。\\

% 「私も頑張らなきゃだし、周達も頑張って。まあ、今の内に羽根を伸ばしておいた方がいいとは思いますけどー？」
「我也必须努力了，周你们也要加油。不过我觉得，还是趁现在先放松一下比较好哦～？」

% 「お前は伸ばし過ぎなの」\\
「你放松过头了」\\

% 　羽根というより足を伸ばして手も大きく伸びをするように伸ばした千歳に呆れてつっこめば、千歳から「メリハリですのでー」と笑い声が帰ってきた。\\
看着浑身上下伸展腿脚，连手都舒展着伸懒腰的千岁，周无奈地吐槽道。千岁笑着回答「这叫劳逸结合」。\\

% 　そのメリハリを是非ともしっかりつけていただきたいものだったが、千歳も遊び呆ける事はないと分かっているし自制するだろうと納得した周は、視線で一度釘を刺しつつも余計な事を言わないよう冷めた紅茶で口を塞いだ。\\
周倒是非常希望她能好好把握这所谓的劳逸结合，但他知道千岁也不会一味地只顾着玩，应该会自我克制的。于是周在用眼神稍微敲打了一下她的同时，为了不让自己说出多余的话，用冷掉的红茶堵住了自己的嘴。\\

% 「春休みにしなければならない事、今しか出来ない事、今やりたい事……ですか」
「春假必须要做的事，只有现在能做的事，现在想做的事……吗」

% 「真昼？」
「真昼？」

% 「いえ、何でもないです」\\
「不，没什么」\\

% 　小さな声で何かを呟いた真昼に聞き返すが、真昼は首を振って静かに曖昧に微笑むだけだった。
真昼小声嘟囔着什么，周反问了一句，但真昼只是摇了摇头，静静地、含糊地微笑着。

